\documentclass[12pt,a4paper,oneside,openany]{book}

% =================================================================
% Pacotes e configuração geral
% (a ordem importa: hyperref SEMPRE por último)
% =================================================================
\usepackage[T1]{fontenc}
\usepackage[utf8]{inputenc}
\usepackage[brazil]{babel}
\usepackage{lmodern}
\usepackage{microtype}
\usepackage[a4paper,top=2cm,bottom=2cm,left=2.5cm,right=2cm]{geometry}
\usepackage[svgnames]{xcolor}
\usepackage{tikz}
\usepackage{graphicx}
\usepackage{booktabs}
\usepackage{array}
\usepackage{enumitem}
\usepackage{fancyhdr}
\usepackage{titlesec}
\usepackage{pdfpages}
\usepackage[noautomatic]{imakeidx}
\usepackage[hidelinks,bookmarksopen=true]{hyperref}

% --- cores (sem definições duplicadas) ---
\definecolor{SASNavy}{HTML}{003B63}
\definecolor{SASBlue}{HTML}{07517D}
\definecolor{SASOrange}{HTML}{F47721}
\definecolor{SASGold}{HTML}{F5A900}
\definecolor{SASLight}{HTML}{F4F7F9}
\definecolor{SASGray}{HTML}{536574}

\hypersetup{
  pdftitle={Anais do III Simpósio de Automação e Sistemas},
  pdfauthor={Laboratório de Automação -- DCA/UFRN},
  pdfsubject={III SAS -- 12 e 13 de novembro de 2026},
  pdfkeywords={automação; sistemas; controle; inteligência artificial; anais}
}

% --- títulos ---
\titleformat{\section}
  {\fontfamily{phv}\selectfont\Large\bfseries\color{DarkGoldenrod}}
  {\thesection}{1em}{}
\titleformat{\subsection}
  {\fontfamily{phv}\selectfont\large\bfseries\color{DarkGoldenrod}}
  {\thesubsection}{1em}{}
\titleformat{\subsubsection}
  {\fontfamily{phv}\selectfont\normalsize\bfseries\color{DarkGoldenrod}}
  {\thesubsubsection}{1em}{}
\titleformat{\chapter}
  {\fontfamily{phv}\selectfont\huge\bfseries\color{DarkGoldenrod}}
  {\thechapter}{1em}{}

% =================================================================
% Cabeçalho/rodapé e NUMERAÇÃO DE PÁGINAS
% O número fica SEMPRE no rodapé central, em todas as páginas
% (inclusive nas páginas dos artigos incluídos via pdfpages).
% =================================================================
\setlength{\headheight}{15pt}


\newcommand{\PaginaDourada}{%
 \begin{tikzpicture}[remember picture,overlay]
    \begin{scope}[shift={(current page.south)}]
      \fill[SASGold] (-0.45,0) [rounded corners=4pt]
        -- (-0.45,0.9) -- (0.45,0.9)
        [sharp corners] -- (0.45,0) -- cycle;
      \node[text=SASNavy, font=\sffamily\footnotesize\bfseries]
        at (0,0.45) {\thepage};
    \end{scope}
  \end{tikzpicture}}


\fancypagestyle{anais}{%
  \fancyhf{}%
  \fancyhead[R]{\small III Simpósio de Automação e Sistemas}%
  \fancyfoot[C]{\small \thepage}%
  \renewcommand{\headrulewidth}{0.4pt}%
  \renewcommand{\footrulewidth}{0pt}}
\fancypagestyle{plain}{%   1ª página de \chapter*: mesmo rodapé, sem filete
  \fancyhf{}%
  \fancyfoot[C]{\small\thepage}%
  \renewcommand{\headrulewidth}{0pt}}
\fancypagestyle{artigo}{%  páginas dos PDFs dos artigos: só o número
  \fancyhf{}%
  \fancyfoot[C]{{\PaginaDourada}}%
  \renewcommand{\headrulewidth}{0pt}}
\pagestyle{anais}

% =================================================================
% Macro para inserir artigos
%   \Artigo{arquivo-sem-extensão}{Título no sumário}{\Autor{..}\Autor{..}}
% - A entrada do sumário é criada DENTRO da 1ª página do artigo (pagecommand),
%   e não depois do \includepdf: isso corrige o deslocamento de um
%   artigo na numeração do sumário.
% - O índice de autores aponta para a 1ª página do artigo.
% =================================================================
\newcommand{\Autor}[1]{\index[autores]{#1}}
\newcommand{\Artigo}[3]{%
  \includepdf[pages=1,
     pagecommand={\thispagestyle{artigo}\phantomsection
                  \addcontentsline{toc}{section}{#2}#3}]{artigos/#1.pdf}%
  \includepdf[pages=2-,
     pagecommand={\thispagestyle{artigo}}]{artigos/#1.pdf}}

\makeindex[name=autores,title={Índice de autores},columns=2,intoc]

\begin{document}

% =================================================================
% Capa
% =================================================================
\begin{titlepage}
  \thispagestyle{empty}
  \begin{tikzpicture}[remember picture,overlay]
    \fill[SASLight] (current page.south west) rectangle (current page.north east);
    \fill[SASNavy] (current page.north west) rectangle
    ([yshift=-0.48cm]current page.north east);
    \fill[SASGold] ([yshift=-0.48cm]current page.north west) rectangle
    ([yshift=-0.56cm]current page.north east);
    \fill[SASOrange] (current page.south west) rectangle
    ([yshift=0.32cm]current page.south east);
  \end{tikzpicture}

  \vspace*{0.55cm}
  \begin{center}
    {\sffamily\bfseries\color{SASNavy}\Large III SIMPÓSIO DE AUTOMAÇÃO E SISTEMAS}\\[0.18cm]

    \vspace{1.0cm}
    % não quebra a compilação se o logo estiver ausente
    \IfFileExists{logo-sas-color.pdf}%
    {\includegraphics[width=0.28\textwidth]{logo-sas-color.pdf}}%
    {\fbox{\rule{0pt}{2.5cm}\hspace{4cm}\textsf{logo}}}

    \vspace{0.9cm}
    {\sffamily\bfseries\color{SASOrange}\large \textbullet\ ANAIS \textbullet\ }\\[1cm]
    {\sffamily\bfseries\color{SASNavy}\fontsize{25}{30}\selectfont
    III Simpósio\\[12pt] de Automação e Sistemas}

    \vspace{0.5cm}
    {\color{SASGray}\rule{0.48\textwidth}{1.1pt}}

    \vspace{0.5cm}
    {\sffamily\color{SASGray}\large Pesquisa, inovação e aplicações em automação}
  \end{center}

  \vfill
  \begin{center}
    \begin{tikzpicture}
      \node[fill=white,draw=SASBlue!25,rounded corners=3pt,
      inner xsep=18pt,inner ysep=12pt,text width=0.72\textwidth,align=center] {
      {\sffamily\bfseries\color{SASNavy} LAUT / DCA / CT / UFRN}\\[0.22cm]
      {\sffamily\color{SASGray} NATAL -- RN}\\[0.18cm]
      {\sffamily\bfseries\color{SASOrange} 2026}
      };
    \end{tikzpicture}

    \vspace{0.65cm}
    {\sffamily\small\color{SASGray} Organizadores: LAUT/DCA/UFRN}
  \end{center}
\end{titlepage}

% =================================================================
% Folha de rosto  (página i: contada, sem número visível)
% NOTA: não usar \begin{titlepage} aqui. Em documentos "oneside" esse
% ambiente executa \setcounter{page}{1} ao final, reiniciando a
% contagem e fazendo a Apresentação voltar a ser "i".
% =================================================================
\frontmatter
\thispagestyle{empty}
\begin{center}
  \vspace*{2cm}
  {\Large\bfseries Anais do III Simpósio de Automação e Sistemas\par}
  \vspace{0.5cm}
  {\large III SAS\par}
  \vfill
  {\large Organização\\Laboratório de Automação (LAUT)\\
    Departamento de Engenharia de Computação e Automação\\
    Universidade Federal do Rio Grande do Norte\par}
  \vfill
  {\large Natal -- RN\\2026\par}
\end{center}
\clearpage

% =================================================================
% Apresentação
% =================================================================
\chapter*{Apresentação}
\addcontentsline{toc}{chapter}{Apresentação}

O III Simpósio de Automação e Sistemas (III SAS) é um evento promovido pelo
Laboratório de Automação (LAUT), vinculado ao Departamento de Engenharia de
Computação e Automação da Universidade Federal do Rio Grande do Norte
(DCA/UFRN). Realizado nos dias 12 e 13 de novembro de 2026, em Natal/RN, o
simpósio reúne atividades de discussão, divulgação científica e apresentação
de trabalhos relacionados à automação e a sistemas inteligentes.

Esta edição também integra as comemorações dos 30 anos do curso de Engenharia
de Computação da UFRN. A programação contempla palestras, sessões plenárias e
apresentações de artigos, abordando temas como inteligência artificial,
desenvolvimento de software, automação e aplicações nos setores produtivos.

\chapter*{Organização e comissões}
\addcontentsline{toc}{chapter}{Organização e comissões}

\section*{Organização}
Laboratório de Automação (LAUT) -- Departamento de Engenharia de Computação e
Automação (DCA/UFRN).

\section*{Comissão Técnica}
\begin{itemize}[itemsep=1pt,topsep=2pt]
  \item Agostinho de Medeiros Brito Junior (UFRN)
  \item Anderson Cavalcanti (UFRN)
  \item André Laurindo Maitelli (UFRN)
  \item Carlos Eduardo Trabuco Dorea (UFRN)
  \item Fábio Meneghetti Ugulino de Araujo (UFRN) % confirmar grafia (acento)
\end{itemize}

\section*{Comissão de Apoio}
\begin{itemize}[itemsep=1pt,topsep=2pt]
  \item Ana Sophia Diógenes Soares (UFRN) % confirmar grafia (acento)
  \item Aurea da Silva Cabral (UFRN)
  \item Dominic N. Rodrigues (UFRN)
  \item Fernando Henrique do Couto (UFRN)
  \item Guilherme C. M. Nascimento (UFRN)
  \item João Paulo Ferreira (UFRN)
  \item Laíze Suélia da Silva Oliveira (UFRN)
  \item Marina Lucena Araujo (UFRN)
\end{itemize}

% =================================================================
% Programação (espaçamento reduzido para não deixar poucos itens
% órfãos numa página isolada)
% =================================================================
\chapter*{Programação}
\addcontentsline{toc}{chapter}{Programação}

\setlist[description]{leftmargin=2.2cm,style=nextline,itemsep=3pt,parsep=0pt,topsep=3pt}

\section*{12 de novembro de 2026 -- Auditório da Reitoria}
\begin{description}
  \item[08h00] Credenciamento.
  \item[09h00] Abertura -- Prof. Anderson Cavalcanti (Chefe do DCA/UFRN) e Profa. Carla Maitelli (Diretora do Centro de Tecnologia/UFRN).
  \item[09h30] Palestra: \textit{30 anos do curso de Engenharia de Computação} -- Prof. Samuel Xavier (DCA/UFRN).
  \item[10h20] Coffee-break.
  \item[10h50] Palestra: \textit{A disrupção no desenvolvimento de software: onde estamos e o que nos espera} -- Eng. Gleydson Lima (CEO, eSIG Group).
  \item[11h40] Palestra: \textit{O curso de Engenharia de Computação na UFRN -- Evolução} -- Prof. Adelardo Medeiros (DCA/UFRN).
  \item[12h30] Almoço.
\end{description}

\section*{12 de novembro de 2026 -- Auditório do CTEC}
\begin{description}
  \item[14h00] Plenária: \textit{Da automação à autonomia: como a inteligência artificial está redefinindo a engenharia} -- Prof. Ivanovitch Silva (DCA/UFRN).
  \item[15h00] Coffee-break.
  \item[15h30] Apresentação de artigos (15 minutos de apresentação e 5 minutos para perguntas) -- 6 trabalhos.
\end{description}

\section*{13 de novembro de 2026 -- Auditório do CTEC}
\begin{description}
  \item[09h00] Plenária: \textit{Automação aplicada ao setor de óleo e gás} -- M.Sc. Heitor Gomes (CEO, Idus Inteligência) e Prof. André Maitelli (DCA/UFRN).
  \item[10h00] Coffee-break.
  \item[10h30] Apresentação de artigos (15 minutos de apresentação e 5 minutos para perguntas) -- 6 trabalhos.
  \item[12h30] Almoço.
  \item[14h00] Apresentação de artigos (15 minutos de apresentação e 5 minutos para perguntas) -- 6 trabalhos.
  \item[16h00] Coffee-break.
  \item[16h30] Encerramento.
\end{description}

% =================================================================
% Sumário e trabalhos
% =================================================================
\tableofcontents
\mainmatter

\chapter*{Trabalhos publicados}
\addcontentsline{toc}{chapter}{Trabalhos publicados}

\Artigo{01_estudos_preliminares_para_o_desenvolvimento_de_uma_arquitetura_modular_baseada_em_ia_para_interacao_com_sistemas_de_gestao_de_malhas_de_controle}
{Estudos preliminares para o desenvolvimento de uma arquitetura modular baseada em IA para interação com sistemas de gestão de malhas de controle}
{\Autor{Ana Sophia Diógenes Soares}\Autor{Fernando Eduardo Dias}\Autor{Fábio Meneghetti Ugulino de Araujo}}

% TODO: preencher autores -> {\Autor{Sobrenome, Nome}\Autor{...}}
\Artigo{02_sistema_de_inspecao_visual_automatizada_para_controle_de_qualidade_em_pecas_fabricadas}
{Sistema de inspeção visual automatizada para controle de qualidade em peças fabricadas}
{\Autor{Israel Lucas}}

% TODO: preencher autores -> {\Autor{Sobrenome, Nome}\Autor{...}}
\Artigo{03_analise_teorica_e_experimental_da_resposta_de_um_gerador_sincrono_a_insercao_de_um_degrau_de_carga}
{Análise teórica e experimental da resposta de um gerador síncrono à inserção de um degrau de carga}
{\Autor{Levi da Silva Lopes}\Autor{José Higor Rodrigues Soares}\Autor{Rodrigo Prado de Medeiros}\Autor{Thiago de Oliveira Alves Rocha}\Autor{Ricardo Lúcio de Araújo Ribeiro}\Autor{José Patrocínio da Silva}}

% TODO: preencher autores -> {\Autor{Sobrenome, Nome}\Autor{...}}
\Artigo{04_analise_da_influencia_de_variacoes_de_atraso_em_controladores_gpc_aplicados_ao_modelo_de_slate_para_pressao_arterial}
{Análise da influência de variações de atraso em controladores GPC aplicados ao modelo de slate para pressão arterial}
{\Autor{Fernando Henrique de Couto}\Autor{Aurea da Silva Cabral}\Autor{André Laurindo Maitelli}}

% TODO: preencher autores -> {\Autor{Sobrenome, Nome}\Autor{...}}
\Artigo{05_gemeos_digitais_na_industria_4_0_aplicacao_em_tanques_acoplados}
{Gêmeos digitais na indústria 4.0: aplicação em tanques acoplados}
{\Autor{Kaillane Araújo da Silva}\Autor{Fabio Meneghetti Ugulino de Araujo}}

% TODO: preencher autores -> {\Autor{Sobrenome, Nome}\Autor{...}}
\Artigo{06_uma_revisao_de_metodologias_inteligentes_de_classificacao_de_curtos_circuitos_em_sistemas_eletricos_de_potencia}
{Uma revisão de metodologias inteligentes de classificação de curtos circuitos em sistemas elétricos de potência}
{\Autor{Guilherme Nascimento da Silva}, \Autor{Fabio Meneghetti Ugulino de Araujo}, \Autor{Rodrigo Prado de Medeiros}}

% TODO: preencher autores -> {\Autor{Sobrenome, Nome}\Autor{...}}
\Artigo{07_interface_baseada_em_agentes_de_ia_para_exploracao_e_controle_de_qualidade_de_dados_sismicos}
{Interface baseada em agentes de IA para exploração e controle de qualidade de dados sísmicos}
{\Autor{José A. A. Oliveira}, \Autor{Luiz F. S. Silva}, \Autor{Ernane Ferreira Rocha Junior}, \Autor{Victor Gomes}, \Autor{Ignacio Sánchez-Gendriz}, \Autor{Ivanovitch Silva}, \Autor{Luiz Affonso Guedes}, \Autor{Tiago Barros}}

% TODO: preencher autores -> {\Autor{Sobrenome, Nome}\Autor{...}}
\Artigo{08_identificacao_de_padroes_e_anomalias_em_reclamacoes_de_telecomunicacoes_no_nordeste_brasileiro_usando_aprendizado_nao_supervisionado}
{Identificação de padrões e anomalias em reclamações de telecomunicações no Nordeste brasileiro usando aprendizado não supervisionado}
{\Autor{Sayonara Cristina de Oliveira Magalhães}, \Autor{Emelyn Clementino Freire}, \Autor{Arthur Phelipe de Oliveira Queiroz}, \Autor{Marcelo Jorge Oliveira de Albuquerque Melo}}

% TODO: preencher autores -> {\Autor{Sobrenome, Nome}\Autor{...}}
\Artigo{10_biblioteca_para_calculo_de_conjuntos_invariantes_com_julia}
{Biblioteca para cálculo de conjuntos invariantes com Julia}
{\Autor{Nicholas W. Lopes}, \Autor{Felipe R. Silva}, \Autor{Carlos E. T. Dórea}}

% TODO: preencher autores -> {\Autor{Sobrenome, Nome}\Autor{...}}
\Artigo{12_desenvolvimento_de_ferramenta_computacional_multicriterio_para_avaliacao_de_desempenho_de_controladores_pid}
{Desenvolvimento de uma ferramenta computacional multicritério para avaliação de desempenho de controladores PID}
{\Autor{Thyfanne Dayane Cosme do Nascimento}, \Autor{Dominic Nascimento Rodrigues}, \Autor{André Laurindo Maitelli}}

% TODO: preencher autores -> {\Autor{Sobrenome, Nome}\Autor{...}}
\Artigo{13_framework_paralelo_de_ajuste_de_campos_de_forca_atomisticos_e_coarse_grain_de_grafeno_atraves_de_otimizacao_por_enxame_de_particulas_pso}
{Framework paralelo de ajuste de campos de força atomísticos e coarse-grain de grafeno através de otimização por enxame de partículas PSO}
{\Autor{Iwri Simplicio Felipe}, \Autor{Elton José Figueiredo de Carvalho}}

% TODO: preencher autores -> {\Autor{Sobrenome, Nome}\Autor{...}}
\Artigo{14_plataforma_linux_embarcada_para_integracao_de_simulador_de_bombeio_mecanico_com_controlador_weg_alc11}
{Plataforma Linux embarcada para integração de simulador de bombeio mecânico com controlador WEG ALC11}
{\Autor{Diego Moura}, \Autor{Andres Ortiz}, \Autor{André Laurindo Maitelli}, \Autor{Rodrigo Dacol}}

% TODO: preencher autores -> {\Autor{Sobrenome, Nome}\Autor{...}}
\Artigo{15_ressintonia_automatica_de_controladores_pid_baseada_em_estimativa_de_robustez}
{Ressintonia automática de controladores PID baseada em estimativa de robustez}
{\Autor{Araújo, Marina L.}, \Autor{Dórea, C. E. Trabuco}}

% TODO: preencher autores -> {\Autor{Sobrenome, Nome}\Autor{...}}
\Artigo{16_industria_sem_telas_perspectivas_para_o_uso_de_dispositivos_vestiveis_assistidos_por_modelos_de_linguagem_de_pequeno_porte}
{Indústria sem telas: perspectivas para o uso de dispositivos vestíveis assistidos por modelos de linguagem de pequeno porte}
{\Autor{Mateus Araujo}, \Autor{Leonardo Cavalcanti}, \Autor{Alice Freire}, \Autor{Morsinaldo Medeiros}, \Autor{Hilton Machado}, \Autor{Erick Justino}, \Autor{Marianne Silva}, \Autor{Carlos M. D. Viegas}, \Autor{Ivanovitch Silva}}

\Artigo{17_arquitetura_edgecloud_para_monitoramento_inteligente_de_veiculos_conectados_com_auditabilidade_via_blockchain}
{Arquitetura edge-cloud para monitoramento inteligente de veículos conectados com auditabilidade via blockchain}
{\Autor{Miguel Euripedes}, \Autor{Morsinaldo Medeiros}, \Autor{Thais Medeiros}, \Autor{Hilton Machado}, \Autor{Erick Justino}, \Autor{Alice Freire}, \Autor{Marianne Silva}, \Autor{Carlos M. D. Viegas}, \Autor{Ivanovitch Silva}}

\Artigo{18_desenvolvimento_de_uma_ferramenta_computacional_em_aplicacao_web_para_simular_o_comportamento_dinamico_do_bombeio_por_cavidades_progressivas_bcp_e_estudar_o_fenomeno_de_reversao_d}
{Desenvolvimento de uma ferramenta computacional em aplicação WEB para simular o comportamento dinâmico do bombeio por cavidades progressivas (BCP) e estudar o fenômeno de reversão do fluxo}
{\Autor{Guilherme da Costa Medeiros Nascimento}, \Autor{André Laurindo Maitelli}, \Autor{Camila Raquel Sena de Almeida}, \Autor{Pedro Manuel Pereira Chamorro}, \Autor{Rutácio de Oliveira Costa} }

% =================================================================
% Elementos pós-textuais
% =================================================================
\backmatter

\printindex[autores]

\chapter*{Expediente}
\addcontentsline{toc}{chapter}{Expediente}

\noindent\textbf{III Simpósio de Automação e Sistemas -- III SAS}\\
12 e 13 de novembro de 2026 -- Natal/RN\\
Organização: Laboratório de Automação (LAUT), DCA/UFRN.

\vspace{0.5cm}
\noindent Anais organizados pela Comissão Técnica do III SAS.\\
\noindent Apoio: FUNPEC, Idus Inteligência, eSIG Group, Centro de Tecnologia, DCA, LAUT.

\end{document}
